\documentclass{beamer}
\hypersetup{colorlinks=true, urlcolor=blue}
\usepackage{listings}
\lstdefinestyle{demo}{
    basicstyle=\ttfamily\small,
    backgroundcolor=\color{black!5},
    frame=single,
    breaklines=true,
    columns=fullflexible,
    keepspaces=true,
}
% Inline code: verbatim argument, so underscores/braces/etc. need no escaping.
% Requires the enclosing frame to be [fragile].
\NewDocumentCommand{\code}{v}{\texttt{#1}}
%Information to be included in the title page:
\title{A Gentle Introduction to Python for Research Computing}
\author{Ben Keene}
\institute{University of Central Florida}
\date{22 June 2026}

\begin{document}

\frame{\titlepage}

\begin{frame}
    \frametitle{License \& Attribution}
    \begin{center}
        This material is adapted from \textbf{The Carpentries}\\
        (Software Carpentry, Data Carpentry, Library Carpentry).

        \bigskip

        Licensed under \textbf{Creative Commons Attribution 4.0 International}\\
        (\href{https://creativecommons.org/licenses/by/4.0/}{CC BY 4.0}).

        \bigskip

        You are free to \emph{share} and \emph{adapt} this material for any purpose,\\
        provided you give appropriate credit.

        \bigskip

        \small \url{https://carpentries.org/}
    \end{center}
\end{frame}

\begin{frame}
    \frametitle{Overview}
    In this beginner-friendly session, we will cover the core Python concepts for research workflows, including data handling, basic analysis, and scripting.
\end{frame}

\begin{frame}
    \frametitle{Why Python for Research?}
    \begin{itemize}
        \item High-level language, easy to learn
        \item Large ecosystem of scientific libraries
        \item Strong community support
        \item Great for automation and reproducibility
    \end{itemize}

    \bigskip

    However, Python can be slow for certain tasks, better to use C or Fortran for scientific computing.

    \bigskip

    In ML/Deep Learning, Python libraries like PyTorch and TensorFlow leverage C++/CUDA for performance-critical operations.
\end{frame}
\begin{frame}
    \frametitle{Characterizing Research Workflows}
    \begin{itemize}
        \item Data generation/collection
        \item Data processing and analysis
        \item Visualization and reporting
    \end{itemize}
\end{frame}



\begin{frame}
    \frametitle{Data Generation \& Collection}
    \textbf{Data Generation}: Simulations, experiments, synthetic data used for training machine learning models.

    \textbf{Data Collection}: Gathering data from various sources, manually downloading datasets, using APIs, or scraping data from the web.

    \begin{itemize}
        \item Tabular data: CSV, Excel, SQL databases
        \item Image data: JPEG, PNG, TIFF
        \item Text data: JSON, XML, plain text
    \end{itemize}
\end{frame}

\begin{frame}
    \frametitle{Data Processing \& Analysis}
    \textbf{Data Processing}: Cleaning, transforming, and preparing data for analysis.

    \textbf{Data Analysis}: Probably your specific research needs.

    \begin{itemize}
        \item Training ML models and evaluating their performance.
        \item Simulating various combinations of chemical reactions.
        \item Evaluating the effectiveness of different approaches to a research question.
    \end{itemize}
\end{frame}

\begin{frame}
    \frametitle{Visualization \& Reporting}
    \textbf{Visualization}: Creating plots, charts, and other visual representations of data.

    \begin{itemize}
        \item Matplotlib, Seaborn, Plotly for plotting
        \item Jupyter Notebooks for interactive data analysis
    \end{itemize}

    \textbf{Reporting}: Saving the output of your analysis in a format that can be distributed.
    Key for reproducibility and sharing results.

    \begin{itemize}
        \item Text files for documentation, logging, saving configurations.
        \item Auto-generate a \LaTeX{} document for publication-quality output.
    \end{itemize}
\end{frame}

\begin{frame}
    \frametitle{Getting Started}
    Navigate to the following URL to get started: \url{https://bit.ly/rcijupyter}.

    \bigskip

    Your username is the first letter of your first name, and your full last name.

    $$\text{John Doe} \mapsto \text{jdoe}$$
    $$\text{Ima Knight} \mapsto \text{iknight}$$

    \bigskip

    \begin{alertblock}{\centering Remember your password!}
        Your password will be set to whatever you enter in the password field.
        Make sure to remember it for future use.
    \end{alertblock}
\end{frame}

\begin{frame}
    \frametitle{Our Scenario: A Miracle Arthritis Cure?}
    Dr. Maverick claims a new drug cures arthritis inflammation
    flare-ups within \textbf{3 weeks} of a patient's first flare-up.

    \bigskip

    After months of asking, we finally have the clinical trial data:
    a single CSV spreadsheet.

    \bigskip

    \textbf{Our job} is to check the claim. To do that, we will:
    \begin{itemize}
        \item Calculate the average inflammation per day across all patients.
        \item Plot the result to discuss and share with colleagues.
    \end{itemize}

    \bigskip

    This is a typical research workflow: \emph{collect} $\rightarrow$
    \emph{analyze} $\rightarrow$ \emph{visualize} $\rightarrow$ \emph{conclude}.
\end{frame}

\begin{frame}[fragile]
    \frametitle{Interpreting the Data}
    The trial followed \textbf{60 patients} over \textbf{40 days}, stored as
    comma-separated values (CSV):
    \begin{itemize}
        \item Each \textbf{row} is one patient.
        \item Each \textbf{column} is one day of the trial.
        \item Each \textbf{value} is the number of inflammation flare-ups
              that patient had on that day.
    \end{itemize}

    \begin{lstlisting}[style=demo]
0,0,1,3,1,2,4,7,8,3,3,3,10,5,7,4,7,...
0,1,2,1,2,1,3,2,2,6,10,11,5,9,4,4,7,...
0,1,1,3,3,2,6,2,5,9,5,7,4,5,4,15,5,...
    \end{lstlisting}

    Row 3, column 7 is \code{6}: patient 3 had 6 flare-ups on day 7.
    Patients start medication after their first flare-up, so we expect
    inflammation to rise, then fall as the drug takes effect.
\end{frame}

\begin{frame}[fragile]
    \frametitle{Live Demonstration}
    Within JupyterLab, open a new terminal. Type \code{python} and press Enter.


\end{frame}

\begin{frame}[fragile]
    \frametitle{Demo Reference}
    \begin{lstlisting}[style=demo]
# shell
python

# Python REPL
>>> 3 + 5 * 4
23
>>> weight_kg = 60
>>> print(weight_kg)
60
    \end{lstlisting}
\end{frame}

\begin{frame}[fragile]
    \frametitle{Data Types}
    \begin{itemize}
        \item \code{int}: Integer values.
        \item \code{float}: Floating-point values.
        \item \code{str}: String values.
        \item And more\ldots
    \end{itemize}
\end{frame}

\begin{frame}[fragile]
    \frametitle{Using Variables}
    \begin{lstlisting}[style=demo]
>>> x = 5
>>> y = 10
>>> z = x + y
>>> print(z)
15
>>> first_name = 'Ben'
>>> last_name = 'Keene'
>>> full_name = first_name + ' ' + last_name
>>> print(full_name)
Ben Keene
    \end{lstlisting}
\end{frame}

\begin{frame}[fragile]
    \frametitle{Printing Output}
    \begin{itemize}
        \item Use the \code{print()} function to display output.
        \item You can print strings, numbers, and other data types.
    \end{itemize}
    \bigskip

    Either: \code{print(your_variable)} or

    \code{print("Your message:", your_variable)}.

    \bigskip

    Fancy formatting using f-strings:

    \code{print(f"Your message: {your_variable}")}.
\end{frame}

\begin{frame}[fragile]
    \frametitle{Functions}
    We've already used the built-in function \code{print()}.

    \bigskip

    \code{type()} returns the type of an object.

    \code{help()} provides help information for a function or module.

    \bigskip

    You can define your own functions using the \code{def} keyword.

    \begin{lstlisting}[style=demo]
>>> def my_function():
>>>     print("Hello, World!")
>>> my_function()
Hello, World!
    \end{lstlisting}
\end{frame}

\begin{frame}[fragile]
    \frametitle{Loading Data}
    Make sure you are inside the \code{beginner-python} directory in your shell.

    \bigskip

    If you are lost, you can enter \code{cd ~/beginner-python}.
\end{frame}

\begin{frame}[fragile]
    \frametitle{Loading Data 2}
    \begin{lstlisting}[style=demo]
>>> import numpy as np
>>> np.loadtxt(fname='data/inflammation-01.csv', delimiter=',')
>>> data = np.loadtxt(fname='data/inflammation-01.csv', delimiter=',')
>>> print(data)
...
>>> print(type(data))
<class 'numpy.ndarray'>
...
>>> print(data.dtype)
float64
>>> print(data.shape)
(60, 40)        # 60 rows 40 columns
    \end{lstlisting}
\end{frame}

\begin{frame}[fragile]
    \frametitle{Visualizing the Data}
    Let's use a Jupyter notebook (ipynb) to visualize the data.

    \bigskip

    From the notebook, we will import Matplotlib and create some plots to investigate the claim that the medicine is effective.
\end{frame}

\begin{frame}[fragile]
    \frametitle{A First Plot: The Heat Map}
    If you started a fresh notebook, reload the data first:
    \begin{lstlisting}[style=demo]
import numpy as np
data = np.loadtxt(fname='data/inflammation-01.csv', delimiter=',')
    \end{lstlisting}

    Now import Matplotlib and draw a heat map of the whole dataset:
    \begin{lstlisting}[style=demo]
import matplotlib.pyplot as plt
image = plt.imshow(data)
plt.show()
    \end{lstlisting}

    Each \textbf{row} is a patient, each \textbf{column} a day.
    Blue is low, yellow is high.
\end{frame}

\begin{frame}[fragile]
    \frametitle{Reading the Heat Map}
    The heat map matches Dr. Maverick's story \emph{so far}:
    \begin{itemize}
        \item Patients start medication once flare-ups begin.
        \item It takes around 3 weeks for the drug to take effect.
        \item Flare-ups appear to drop to zero by the end of the trial.
    \end{itemize}

    \bigskip

    But a single picture can hide a lot. Let's look at some summary
    statistics \emph{across all patients}, one value per day.
\end{frame}

\begin{frame}[fragile]
    \frametitle{Average Inflammation Over Time}
    \code{np.mean(data, axis=0)} collapses the 60 rows into a single
    average for each of the 40 days:
    \begin{lstlisting}[style=demo]
ave_inflammation = np.mean(data, axis=0)
ave_plot = plt.plot(ave_inflammation)
plt.show()
    \end{lstlisting}

    A reasonably linear rise and fall --- consistent with the claim that
    the medication takes about 3 weeks to work.

    \bigskip

    A good data scientist never stops at the average.
\end{frame}

\begin{frame}[fragile]
    \frametitle{Maximum and Minimum}
    \begin{lstlisting}[style=demo]
max_plot = plt.plot(np.max(data, axis=0))
plt.show()

min_plot = plt.plot(np.min(data, axis=0))
plt.show()
    \end{lstlisting}

    \begin{itemize}
        \item The \textbf{maximum} rises and falls in clean straight lines.
        \item The \textbf{minimum} looks like a step function.
    \end{itemize}

    Neither trend is biologically plausible. Either our calculation is
    wrong, or \emph{something is wrong with the data}.
\end{frame}

\begin{frame}[fragile]
    \frametitle{Grouping Plots: Subplots}
    We can place several plots in one figure to compare them at a glance.
    \begin{lstlisting}[style=demo]
fig = plt.figure(figsize=(10.0, 3.0))
axes1 = fig.add_subplot(1, 3, 1)
axes2 = fig.add_subplot(1, 3, 2)
axes3 = fig.add_subplot(1, 3, 3)

axes1.set_ylabel('average')
axes1.plot(np.mean(data, axis=0))
axes2.set_ylabel('max')
axes2.plot(np.max(data, axis=0))
axes3.set_ylabel('min')
axes3.plot(np.min(data, axis=0))

fig.tight_layout()
plt.savefig('inflammation.png')
plt.show()
    \end{lstlisting}
\end{frame}

\begin{frame}[fragile]
    \frametitle{Understanding \texttt{add\_subplot}}
    \code{plt.figure(figsize=(w, h))} creates the canvas.

    \bigskip

    \code{fig.add_subplot(rows, cols, n)} adds one plot:
    \begin{itemize}
        \item \textbf{rows}: total rows of subplots
        \item \textbf{cols}: total columns of subplots
        \item \textbf{n}: which cell, counting left-to-right, top-to-bottom
    \end{itemize}
\end{frame}

\begin{frame}
    \frametitle{Optional: Lists}
    Walk through lists, slicing, and other list operations.
\end{frame}

\begin{frame}
    \frametitle{Optional: Loops}
    Walk through \texttt{for} loops and how to use them with lists.
\end{frame}

\begin{frame}[fragile]
    \frametitle{Analyzing Many Files at Once}
    Dr. Maverick eventually hands over \textbf{twelve} CSV files, one per
    clinical trial. We don't want to copy-paste our analysis twelve times.

    \bigskip

    Instead, we'll:
    \begin{itemize}
        \item Get a list of every file whose name matches a pattern.
        \item Write a \code{for} loop to run the same analysis on each one.
    \end{itemize}

    \bigskip

    This is the payoff of writing reusable code: the work scales from one
    file to many for free.
\end{frame}

\begin{frame}[fragile]
    \frametitle{Finding Files with \texttt{glob}}
    The \code{glob} library finds files whose names match a pattern:
    \begin{itemize}
        \item \code{*} matches zero or more characters.
        \item \code{?} matches any single character.
    \end{itemize}

    \begin{lstlisting}[style=demo]
import glob
print(glob.glob('data/inflammation*.csv'))
['data/inflammation-05.csv', 'data/inflammation-11.csv',
 'data/inflammation-12.csv', 'data/inflammation-08.csv',
 ...
 'data/inflammation-01.csv']
    \end{lstlisting}

    \code{glob.glob} returns a \textbf{list} in arbitrary order, so we can
    loop over it.
\end{frame}

\begin{frame}[fragile]
    \frametitle{Sorting the File List}
    The order is arbitrary, so we use \code{sorted()} to get a predictable,
    alphabetical list. We'll start with just the first three files:
    \begin{lstlisting}[style=demo]
import glob
import numpy as np
import matplotlib.pyplot as plt

filenames = sorted(glob.glob('data/inflammation*.csv'))
filenames = filenames[0:3]
    \end{lstlisting}

    \code{sorted()} returns a \emph{new} sorted list; slicing
    \code{[0:3]} keeps the first three entries.
\end{frame}

\begin{frame}[fragile]
    \frametitle{One Loop, Many Files}
    \begin{lstlisting}[style=demo]
for filename in filenames:
    print(filename)
    data = np.loadtxt(fname=filename, delimiter=',')

    fig = plt.figure(figsize=(10.0, 3.0))
    axes1 = fig.add_subplot(1, 3, 1)
    axes2 = fig.add_subplot(1, 3, 2)
    axes3 = fig.add_subplot(1, 3, 3)

    axes1.set_ylabel('average')
    axes1.plot(np.mean(data, axis=0))
    axes2.set_ylabel('max')
    axes2.plot(np.max(data, axis=0))
    axes3.set_ylabel('min')
    axes3.plot(np.min(data, axis=0))

    fig.tight_layout()
    plt.show()
    \end{lstlisting}
\end{frame}

\begin{frame}[fragile]
    \frametitle{Comparing the Results}
    The loop produces one figure per file. Looking across them:
    \begin{itemize}
        \item Files 1 and 2 look almost identical: the same suspiciously
              clean maxima and minima we saw before.
        \item File 3 has much \emph{noisier} average and maximum plots ---
              actually \emph{less} suspicious.
        \item But file 3's minimum is stuck at \textbf{zero} for every
              single day.
    \end{itemize}

    \bigskip

    A heat map of file 3 shows zeros scattered throughout, and one patient
    with \emph{no} flare-ups at all --- they may not even have arthritis.
\end{frame}

\begin{frame}[fragile]
    \frametitle{The Verdict}
    Investigating all twelve files, two categories emerge:
    \begin{itemize}
        \item ``Ideal'' datasets that match the claim perfectly --- but have
              impossible, too-clean maxima and minima.
        \item ``Noisy'' datasets with sporadic missing values and
              unsuitable participants.
    \end{itemize}

    \bigskip

    Worse, three of the noisy files (\code{03}, \code{08}, \code{11}) are
    \textbf{identical down to the last value}.

    \bigskip

    Confronted, Dr. Maverick admits the data was \textbf{fabricated} ---
    fake datasets, padded out with duplicate copies of a bad one.
\end{frame}


\end{document}